\paragraph{生産者物価テイラールール（TR-PPI）}
\label{sec:chapter5_beta_shock_policies_tr_ppi}
式
\eqref{form:chapter3_policies_monetary_home_tr_ppi_i_H_notional_eq}および
\eqref{form:chapter3_policies_monetary_home_tr_ppi_i_H_eq}
により記述されたとおり、TR-PPIの式は以下のようであった。
\begin{align}
\label{form:chapter5_beta_shock_policies_tr_ppi_i_H_notional_eq}
    i_t^{H,notional}&=\rho_i^Hi_{t-1}^{H,notional}+(1-\rho_i^H)\left(i_{ss}^H+\phi_{\pi}^H\log(\pi_t^H)+\phi_y^H(\log(y_t^H)-\log(y_t^{H,pot}))\right)+\varepsilon_t^{i,H} \\
\label{form:chapter5_beta_shock_policies_tr_ppi_i_H_eq}
    i_t^H&=i_t^{H,notional}
\end{align}
ここで
\(\phi_{\pi}^H, \phi_y^H\)は反応係数、
\(\rho_i^H\)は金利平滑化係数、
\(\pi_t^H\)は生産者物価インフレ率、
\(y_t^H\)は実質産出量、
\(y_t^{H,pot}\)は潜在産出量、
\(\varepsilon_t^{i,H}\)は外生ショック項である。
\begin{enumerate}
    \item
    \label{itm:chapter5_beta_shock_policies_tr_ppi_c_H_W}
        \(c_t^{H\to W}\)：
        総消費指数\(c_t^{H\to W}\)の図
        \ref{fig:chapter5_beta_shock_graphs_c_H_W}
        を見ると、TR-PPIは消費者物価テイラールール（TR-CPI）や消費者物価水準目標（CPLT）に見られるような
        極端な大暴落を免れていることがわかる。
        100期までの前半において、落ち込みの幅は
        総消費水準目標（CLT）や名目総消費水準目標（NCLT）には及ばないものの、ある程度の水準を維持している。
        消費の低下は厚生を悪化させる方向に働くため、
        この前半における厚生の押し下げはTR-CPIなどよりも小さく防がれている。
        一方で中盤以降のプラス圏での推移を見ると、厚生を改善させる消費の増加幅は比較的小規模にとどまっている。TR-CPIが遅れて形成する巨大な山や、NCLTが長期にわたって維持する高い水準と比較すると、後半の消費増加による厚生の押し上げ効果は限定的となっている。
    \item
    \label{itm:chapter5_beta_shock_policies_tr_ppi_y_H}
        \(y_t^H\)：
        生産\(y_t^H\)の図
        \ref{fig:chapter5_beta_shock_graphs_y_H}
        を見ると、TR-PPIの軌道は総消費指数\(c_t^{H\to W}\)と同様に、
        消費者物価テイラールール（TR-CPI）や消費者物価水準目標（CPLT）のような
        極端な変動を免れていることがわかる。
        100期までの前半において、TR-PPIの落ち込みは
        総消費水準目標（CLT）や名目総消費水準目標（NCLT）よりは大きいが、
        TR-CPIやCPLTのような極端な落ち込みには至っていない。
        生産の低下は労働の減少を意味するため厚生にはプラスに働くが、
        TR-PPIにおける前半の厚生押し上げ効果はTR-CPIやCPLTほど大きくはない。
        一方で中盤以降のプラス圏での推移（オーバーシュート）を見ると、
        生産の増加は労働の非効用を通じて厚生を悪化させる要因となる。
        TR-PPIのオーバーシュートはCLTやNCLTよりも低く抑えられており、
        またTR-CPIが遅れて形成する巨大な山を作ることもないため、
        後半の生産過剰に伴う厚生のマイナス効果は比較的小さくとどまっている。
    \item
    \label{itm:chapter5_beta_shock_policies_tr_ppi_pi_H}
        \(\pi_t^H\)：
        グロス・インフレ率\(\pi_t^H\)の図
        \ref{fig:chapter5_beta_shock_graphs_pi_H}
        を見ると、TR-PPIの定常状態からの乖離は、
        消費者物価テイラールール（TR-CPI）に見られるような極端な落ち込みや大きな波打ちを免れていることがわかる。
        ショック直後の前半においてもマイナス（デフレ）方向への振れ幅は小さく抑えられており、
        中盤以降のプラス圏（インフレ）への反転も小規模にとどまっている。
        生産者物価インフレ率に反応するテイラールールとして機能しているため、
        TR-CPIと比較してインフレ率のブレが適切に制御されており、
        この定常状態からの乖離の小ささが次項で述べる価格分散の蓄積を抑え、厚生の悪化を防ぐ要因となっている。
    \item
    \label{itm:chapter5_beta_shock_policies_tr_ppi_Delta_t_minus_1_H}
        \(\Delta_{t-1}^H\)：
        価格分散\(\Delta_{t-1}^H\)の図
        \ref{fig:chapter5_beta_shock_welfare_Delta_H}を見ると、
        TR-PPIの定常状態からの乖離は、
        消費者物価テイラールール（TR-CPI）に見られるような中盤における大きな山の形成を免れ、
        極めて低水準に抑えられていることがわかる。
        総消費水準目標（CLT）のような著しい増大もなく、
        生産者物価水準目標（PPLT）や名目GDP水準目標（NGDPLT）と同様に
        全期間を通じてほぼゼロに近い水準を維持している。
        これは前述のインフレ率の変動が適切に制御されていることで価格分散の蓄積が未然に防がれている結果であり、
        この生産効率の悪化の回避がTR-PPIの厚生の低下を防ぐ重要な要因となっている。
\end{enumerate}